\section{Other formalisms and where their encoding costs move}
\label{sec:formalism-encoding-comparison}

The strongest route in this delivery is to retain K230's calibrated
proof enumerator and improve its implementation.  A shorter axiom or a
smaller theory packet is useful only if it reduces the complete blank-tape
machine after the interpreter, finite-proof search and initialization
have been combined.  The alternatives below identify mathematically
sound interfaces for further work while keeping measured packet sizes
separate from measured machine states.

\paragraph{K and the inherited native enumerator.}
K230 already connects its exact nonzero search seed, numerical formula
codes, finite derivations and binary table to ordinary ZFC refutation.
The calibration is unusually concrete: the K-to-ZFC comparisons transform
finite proofs, and the final Rocq theorem names the literal machine.
The main construction preserves this logical interface.  The four
transfer-plus-one normalizations and later source layouts change how
the same search is executed; direct whole-machine simulation establishes
the link to the changed table.  This gives a complete state count and
a usable correctness boundary.\cite{k230}

The shorter K schema retains native first-order equality, irreflexivity
in its encoded presentation, all eight inherited logical roots, and the
four formula tags.  Removing equality rules or switching to a smaller
tag alphabet would require a new completeness and decoding proof.
Source 14's earlier equality-free numerical proposals are a concrete
reason to insist on that obligation: Boolean valuations validate their
emitted rules while making the target false, even though the intended
ordinary set-theoretic schema has the desired consistency strength.
The intended axioms and the actually enumerated calculus must agree.
\footnote{Pinned primary project record:
\url{https://github.com/VladimirReshetnikov/ProveIt/blob/ce192939372ac7944fe2a0208c68a96d49fb5e3b/SetTheory/Cardinals/docs/reports/set-theory/zfc-closure-kernels/intake-146-source-14.md}.}

\paragraph{Direct ZFC and finite arithmetic presentations.}
A direct checker can enumerate finite strings and halt upon finding a
ZFC refutation.  Its calibration is immediate from the proof predicate,
but its table must implement axiom-schema recognition, substitution,
inference checking and fair enumeration.  A finite axiom list can move
these costs into one complicated formula.  For example, for a fixed
$\Pi^0_1$ sentence $\theta$, ordinary ZFC proves
\[
 \Con(\mathsf Q+\theta)\longleftrightarrow\theta.
\]
If $\theta$ fails, its finite witness yields a Robinson-arithmetic proof
of $\neg\theta$; if $\theta$ holds, the natural-number interpretation
and finite-proof soundness give consistency.  Taking
$\theta=\Con(\mathrm{ZFC})$ therefore gives a finite equiconsistent
arithmetic presentation.  Its added sentence still contains the
arithmetized ZFC proof predicate.  This observation supplies no smaller
implemented checker or state count.

\paragraph{Raw closure and finite class-theory routes.}
The inspected empty-input EPSC result removes PSC's distinguished seed
parameter in an equality-free membership language.  The raw schema's
active inputs are all empty objects; support first supplies an empty
object, and raw power, separation and subset-bound constructions recover
the seeded closure scheme.  The dossier gives finite schema-instance
derivations and the inherited ordinary-ZFC consistency comparison.
This is a substantive logical simplification, while its compressed
numerical frontend still needs the same end-to-end calibration as K.

A finite class-theory presentation could replace recognition of arbitrary
schema instances by a fixed list of closure operations.  The specific
18-equation Df3/finite-class-theory route in the inspected dossier stops
at a proposal: it supplies neither the actual ZFC target term nor its
translation.  At that artifact boundary there is no complete calibrated
payload to compare with the delivered machine.  Constructing those two
objects, then measuring their joint interpreter and initializer, would
make the proposal assessable.

\paragraph{E639: bind the descriptor inside its own class.}
The equational supplement gives a more concrete alternative: seven
equations, nine operation symbols and a ground query, encoded in
\textbf{639 bits}.  Relative to the inherited E663 packet, it replaces
\[
 R_b(x,y)=p_\varepsilon(P(y,P(x,b)),V(b,d))
 \quad\hbox{by}\quad
 \widetilde R_b(x,y)=p_\varepsilon(P(y,P(x,b)),b).
\]
There are three occurrences.  Removing two term nodes at each occurrence
saves exactly $3\cdot2\cdot4=24$ bits in the unchanged codec.

The saving needs a representation proof: the descriptor must carry
parameters and simultaneously denote the desired formula class.
With $P$ acting as NAND under class evaluation, put
\[
 T=P(d,P(d,d)),\quad S_{k+1}=T,\quad
 S_i=P(p_i,P(p_i,S_{i+1})),\quad b=P(P(c,c),S_1).
\]
Every $S_i$ evaluates to truth and $b$ evaluates to the same class as $c$.
In the environment $P(y,P(x,b))$, output and input occupy coordinates
zero and one; parameters occupy coordinates $3,5,7,\ldots$.
Build $c$ first using these fixed indices, then build $b$.  Thus parameter
binding is finite and noncircular.  The accompanying proof establishes
both model directions with explicit well-founded constructor weights.
The consistency target is the equations together with $Q(d)\ne d$;
equations alone always admit a one-element model.  This is a proved
24-bit descriptor-binding improvement, with no claimed TM state bound.

\paragraph{What a new equational interpreter must achieve.}
Ground congruence certificates can avoid a general first-order
capture-avoiding substitution engine.  They still require arbitrary
ground instances, either rewrite direction, replacement inside arbitrary
contexts, and fair coverage of every finite derivation.  The earlier
measured E7 generator compiled to 299 working states before supplying
its ZFC-specific query and writer; its isolated emitter used 171.
Those are costs of that particular architecture, neither a compilation
nor a lower bound for E639.\footnote{Pinned primary experiment:
\url{https://github.com/VladimirReshetnikov/ProveIt/blob/13ec39ed43148b1807f51ed6e6528e1c84659562/Computability/BusyBeaver/Research/ZFC/equational_generator_cost.md}.}
The useful next experiment is therefore a complete alternative interpreter
with the 639-bit payload charged, followed by the same source, startup and
unbounded runtime checks used for the K transplants.
